
\documentclass{article}
\usepackage{colortbl}
\usepackage{makecell}
\usepackage{multirow}
\usepackage{supertabular}

\begin{document}

\newcounter{utterance}

\centering \large Interaction Transcript for game `wordle', experiment `high\_frequency\_words\_no\_clue\_no\_critic', episode 4 with gpt{-}5.4{-}azure.
\vspace{24pt}

{ \footnotesize  \setcounter{utterance}{1}
\setlength{\tabcolsep}{0pt}
\begin{supertabular}{c@{$\;$}|p{.15\linewidth}@{}p{.15\linewidth}p{.15\linewidth}p{.15\linewidth}p{.15\linewidth}p{.15\linewidth}}
    \# & \multicolumn{2}{c}{Player} && \multicolumn{2}{c}{Game Master} \\
    \hline

    \theutterance \stepcounter{utterance}  
    & & & \multicolumn{4}{p{0.6\linewidth}}{
        \cellcolor[rgb]{0.9,0.9,0.9}{
            \makecell[{{p{\linewidth}}}]{
                \texttt{\tiny{[P1$\langle$GM]}}
                \texttt{Je bent een taaltovenaar die graag woorden raadt met behulp van de gegeven regels.} \\
\\ 
\texttt{Welkom bij Wordle! Je hebt zes pogingen om het doelwoord te raden, een geldig woord van vijf kleine letters. Gebruik de tags "uitleg:" en "gok:" om voor elke gok een beknopte uitleg te geven.} \\
\\ 
\texttt{Bijvoorbeeld, als je gok "staat" is, moet je antwoord zijn} \\
\texttt{uitleg: dit is een veelvoorkomend woord van vijf letters, en ik begin mijn gok met dit woord.} \\
\texttt{gok: staat} \\
\\ 
\texttt{Na elke gok wordt je antwoord gevalideerd en ontvang je feedback die aangeeft welke letters correct zijn (groen), welke letters correct zijn maar op de verkeerde positie staan (geel), en welke letters onjuist zijn (rood). Deze feedback kan nuttig zijn om te bepalen welke letters je in je volgende gok moet opnemen of juist moet uitsluiten.} \\
\\ 
\texttt{Bijvoorbeeld, de feedback voor "staat" kan zijn:} \\
\texttt{gok\_feedback: s<geel> t<geel> a<groen> a<geel> t<rood>} \\
\\ 
\texttt{De uitleg moet details bevatten over hoe de gok\_feedback wordt gebruikt om tot een nieuwe gok te komen.} \\
\\ 
\texttt{Genereer geen andere tekst.} \\
\\ 
\texttt{Laten we beginnen met je eerste gok.} \\
            }
        }
    }
     \\ \\

    \theutterance \stepcounter{utterance}  
    & \multicolumn{4}{p{0.6\linewidth}}{
        \cellcolor[rgb]{0.561,0.737,0.561}{
            \makecell[{{p{\linewidth}}}]{
                \texttt{\tiny{[P1$\rangle$GM]}}
                \texttt{uitleg: dit is een veelvoorkomend woord van vijf letters met veel gebruikte klinkers en medeklinkers, zodat de eerste feedback veel informatie geeft.} \\
\texttt{gok: stoel} \\
            }
        }
    }
    & & \\ \\

    \theutterance \stepcounter{utterance}  
    & & & \multicolumn{2}{p{0.3\linewidth}}{
        \cellcolor[rgb]{0.9,0.9,0.9}{
            \makecell[{{p{\linewidth}}}]{
                \texttt{\tiny{[GM$|$GM]}}
                \texttt{attempts = 1} \\
\texttt{target\_word = kaart} \\
\texttt{guess = stoel} \\
\texttt{guess\_feedback = s<red> t<yellow> o<red> e<red> l<red>} \\
            }
        }
    }
    & & \\ \\

    \theutterance \stepcounter{utterance}  
    & & & \multicolumn{4}{p{0.6\linewidth}}{
        \cellcolor[rgb]{0.9,0.9,0.9}{
            \makecell[{{p{\linewidth}}}]{
                \texttt{\tiny{[P1$\langle$GM]}}
                \texttt{gok\_feedback: s<rood> t<geel> o<rood> e<rood> l<rood>} \\
\\ 
\texttt{Geef uw antwoord alleen in dit formaat.} \\
\texttt{uitleg: details} \\
\texttt{gok: woord} \\
            }
        }
    }
     \\ \\

    \theutterance \stepcounter{utterance}  
    & \multicolumn{4}{p{0.6\linewidth}}{
        \cellcolor[rgb]{0.561,0.737,0.561}{
            \makecell[{{p{\linewidth}}}]{
                \texttt{\tiny{[P1$\rangle$GM]}}
                \texttt{uitleg: uit de feedback blijkt dat de t wel in het woord zit, maar niet op de tweede positie. De letters s, o, e en l vallen af, dus ik kies een nieuw woord met t op een andere plek en met vooral nieuwe, veelvoorkomende letters.} \\
\texttt{gok: dacht} \\
            }
        }
    }
    & & \\ \\

    \theutterance \stepcounter{utterance}  
    & & & \multicolumn{2}{p{0.3\linewidth}}{
        \cellcolor[rgb]{0.9,0.9,0.9}{
            \makecell[{{p{\linewidth}}}]{
                \texttt{\tiny{[GM$|$GM]}}
                \texttt{attempts = 2} \\
\texttt{target\_word = kaart} \\
\texttt{guess = dacht} \\
\texttt{guess\_feedback = d<red> a<green> c<red> h<red> t<green>} \\
            }
        }
    }
    & & \\ \\

    \theutterance \stepcounter{utterance}  
    & & & \multicolumn{4}{p{0.6\linewidth}}{
        \cellcolor[rgb]{0.9,0.9,0.9}{
            \makecell[{{p{\linewidth}}}]{
                \texttt{\tiny{[P1$\langle$GM]}}
                \texttt{gok\_feedback: d<rood> a<groen> c<rood> h<rood> t<groen>} \\
\\ 
\texttt{Geef uw antwoord alleen in dit formaat.} \\
\texttt{uitleg: details} \\
\texttt{gok: woord} \\
            }
        }
    }
     \\ \\

    \theutterance \stepcounter{utterance}  
    & \multicolumn{4}{p{0.6\linewidth}}{
        \cellcolor[rgb]{0.561,0.737,0.561}{
            \makecell[{{p{\linewidth}}}]{
                \texttt{\tiny{[P1$\rangle$GM]}}
                \texttt{uitleg: nu weten we dat a op de tweede plek staat en t op de vijfde. De letters s, o, e, l, d, c en h zijn uitgesloten, dus ik kies een passend woord in het patroon \_a\_\_t met nieuwe, veelvoorkomende letters om extra informatie te krijgen.} \\
\texttt{gok: markt} \\
            }
        }
    }
    & & \\ \\

    \theutterance \stepcounter{utterance}  
    & & & \multicolumn{2}{p{0.3\linewidth}}{
        \cellcolor[rgb]{0.9,0.9,0.9}{
            \makecell[{{p{\linewidth}}}]{
                \texttt{\tiny{[GM$|$GM]}}
                \texttt{attempts = 3} \\
\texttt{target\_word = kaart} \\
\texttt{guess = markt} \\
\texttt{guess\_feedback = m<red> a<green> r<yellow> k<yellow> t<green>} \\
            }
        }
    }
    & & \\ \\

    \theutterance \stepcounter{utterance}  
    & & & \multicolumn{4}{p{0.6\linewidth}}{
        \cellcolor[rgb]{0.9,0.9,0.9}{
            \makecell[{{p{\linewidth}}}]{
                \texttt{\tiny{[P1$\langle$GM]}}
                \texttt{gok\_feedback: m<rood> a<groen> r<geel> k<geel> t<groen>} \\
\\ 
\texttt{Geef uw antwoord alleen in dit formaat.} \\
\texttt{uitleg: details} \\
\texttt{gok: woord} \\
            }
        }
    }
     \\ \\

    \theutterance \stepcounter{utterance}  
    & \multicolumn{4}{p{0.6\linewidth}}{
        \cellcolor[rgb]{0.561,0.737,0.561}{
            \makecell[{{p{\linewidth}}}]{
                \texttt{\tiny{[P1$\rangle$GM]}}
                \texttt{uitleg: de feedback past precies op het patroon k a a r t: a staat vast op plek twee, t op plek vijf, r zit wel in het woord maar niet op plek drie en k zit wel in het woord maar niet op plek vier. De uitgesloten letters komen hierin niet voor.} \\
\texttt{gok: kaart} \\
            }
        }
    }
    & & \\ \\

    \theutterance \stepcounter{utterance}  
    & & & \multicolumn{2}{p{0.3\linewidth}}{
        \cellcolor[rgb]{0.9,0.9,0.9}{
            \makecell[{{p{\linewidth}}}]{
                \texttt{\tiny{[GM$|$GM]}}
                \texttt{attempts = 4} \\
\texttt{target\_word = kaart} \\
\texttt{guess = kaart} \\
\texttt{guess\_feedback = k<green> a<green> a<green> r<green> t<green>} \\
            }
        }
    }
    & & \\ \\

    \theutterance \stepcounter{utterance}  
    & & & \multicolumn{2}{p{0.3\linewidth}}{
        \cellcolor[rgb]{0.9,0.9,0.9}{
            \makecell[{{p{\linewidth}}}]{
                \texttt{\tiny{[GM$|$GM]}}
                \texttt{game\_result = WIN} \\
            }
        }
    }
    & & \\ \\

\end{supertabular}
}

\end{document}
